% Options for packages loaded elsewhere
\PassOptionsToPackage{unicode}{hyperref}
\PassOptionsToPackage{hyphens}{url}
\PassOptionsToPackage{dvipsnames,svgnames,x11names}{xcolor}
\documentclass[
  11pt,
  letterpaper,
]{article}
\usepackage{xcolor}
\usepackage[margin=1in]{geometry}
\usepackage{amsmath,amssymb}
\setcounter{secnumdepth}{-\maxdimen} % remove section numbering
\usepackage{iftex}
\ifPDFTeX
  \usepackage[T1]{fontenc}
  \usepackage[utf8]{inputenc}
  \usepackage{textcomp} % provide euro and other symbols
\else % if luatex or xetex
  \usepackage{unicode-math} % this also loads fontspec
  \defaultfontfeatures{Scale=MatchLowercase}
  \defaultfontfeatures[\rmfamily]{Ligatures=TeX,Scale=1}
\fi
\usepackage{lmodern}
\ifPDFTeX\else
  % xetex/luatex font selection
\fi
% Use upquote if available, for straight quotes in verbatim environments
\IfFileExists{upquote.sty}{\usepackage{upquote}}{}
\IfFileExists{microtype.sty}{% use microtype if available
  \usepackage[]{microtype}
  \UseMicrotypeSet[protrusion]{basicmath} % disable protrusion for tt fonts
}{}
\makeatletter
\@ifundefined{KOMAClassName}{% if non-KOMA class
  \IfFileExists{parskip.sty}{%
    \usepackage{parskip}
  }{% else
    \setlength{\parindent}{0pt}
    \setlength{\parskip}{6pt plus 2pt minus 1pt}}
}{% if KOMA class
  \KOMAoptions{parskip=half}}
\makeatother
\usepackage{color}
\usepackage{fancyvrb}
\newcommand{\VerbBar}{|}
\newcommand{\VERB}{\Verb[commandchars=\\\{\}]}
\DefineVerbatimEnvironment{Highlighting}{Verbatim}{commandchars=\\\{\}}
% Add ',fontsize=\small' for more characters per line
\newenvironment{Shaded}{}{}
\newcommand{\AlertTok}[1]{\textcolor[rgb]{1.00,0.00,0.00}{\textbf{#1}}}
\newcommand{\AnnotationTok}[1]{\textcolor[rgb]{0.38,0.63,0.69}{\textbf{\textit{#1}}}}
\newcommand{\AttributeTok}[1]{\textcolor[rgb]{0.49,0.56,0.16}{#1}}
\newcommand{\BaseNTok}[1]{\textcolor[rgb]{0.25,0.63,0.44}{#1}}
\newcommand{\BuiltInTok}[1]{\textcolor[rgb]{0.00,0.50,0.00}{#1}}
\newcommand{\CharTok}[1]{\textcolor[rgb]{0.25,0.44,0.63}{#1}}
\newcommand{\CommentTok}[1]{\textcolor[rgb]{0.38,0.63,0.69}{\textit{#1}}}
\newcommand{\CommentVarTok}[1]{\textcolor[rgb]{0.38,0.63,0.69}{\textbf{\textit{#1}}}}
\newcommand{\ConstantTok}[1]{\textcolor[rgb]{0.53,0.00,0.00}{#1}}
\newcommand{\ControlFlowTok}[1]{\textcolor[rgb]{0.00,0.44,0.13}{\textbf{#1}}}
\newcommand{\DataTypeTok}[1]{\textcolor[rgb]{0.56,0.13,0.00}{#1}}
\newcommand{\DecValTok}[1]{\textcolor[rgb]{0.25,0.63,0.44}{#1}}
\newcommand{\DocumentationTok}[1]{\textcolor[rgb]{0.73,0.13,0.13}{\textit{#1}}}
\newcommand{\ErrorTok}[1]{\textcolor[rgb]{1.00,0.00,0.00}{\textbf{#1}}}
\newcommand{\ExtensionTok}[1]{#1}
\newcommand{\FloatTok}[1]{\textcolor[rgb]{0.25,0.63,0.44}{#1}}
\newcommand{\FunctionTok}[1]{\textcolor[rgb]{0.02,0.16,0.49}{#1}}
\newcommand{\ImportTok}[1]{\textcolor[rgb]{0.00,0.50,0.00}{\textbf{#1}}}
\newcommand{\InformationTok}[1]{\textcolor[rgb]{0.38,0.63,0.69}{\textbf{\textit{#1}}}}
\newcommand{\KeywordTok}[1]{\textcolor[rgb]{0.00,0.44,0.13}{\textbf{#1}}}
\newcommand{\NormalTok}[1]{#1}
\newcommand{\OperatorTok}[1]{\textcolor[rgb]{0.40,0.40,0.40}{#1}}
\newcommand{\OtherTok}[1]{\textcolor[rgb]{0.00,0.44,0.13}{#1}}
\newcommand{\PreprocessorTok}[1]{\textcolor[rgb]{0.74,0.48,0.00}{#1}}
\newcommand{\RegionMarkerTok}[1]{#1}
\newcommand{\SpecialCharTok}[1]{\textcolor[rgb]{0.25,0.44,0.63}{#1}}
\newcommand{\SpecialStringTok}[1]{\textcolor[rgb]{0.73,0.40,0.53}{#1}}
\newcommand{\StringTok}[1]{\textcolor[rgb]{0.25,0.44,0.63}{#1}}
\newcommand{\VariableTok}[1]{\textcolor[rgb]{0.10,0.09,0.49}{#1}}
\newcommand{\VerbatimStringTok}[1]{\textcolor[rgb]{0.25,0.44,0.63}{#1}}
\newcommand{\WarningTok}[1]{\textcolor[rgb]{0.38,0.63,0.69}{\textbf{\textit{#1}}}}
\usepackage{longtable,booktabs,array}
\newcounter{none} % for unnumbered tables
\usepackage{calc} % for calculating minipage widths
% Correct order of tables after \paragraph or \subparagraph
\usepackage{etoolbox}
\makeatletter
\patchcmd\longtable{\par}{\if@noskipsec\mbox{}\fi\par}{}{}
\makeatother
% Allow footnotes in longtable head/foot
\IfFileExists{footnotehyper.sty}{\usepackage{footnotehyper}}{\usepackage{footnote}}
\makesavenoteenv{longtable}
\setlength{\emergencystretch}{3em} % prevent overfull lines
\providecommand{\tightlist}{%
  \setlength{\itemsep}{0pt}\setlength{\parskip}{0pt}}
\usepackage{bookmark}
\IfFileExists{xurl.sty}{\usepackage{xurl}}{} % add URL line breaks if available
\urlstyle{same}
\hypersetup{
  colorlinks=true,
  linkcolor={blue},
  filecolor={Maroon},
  citecolor={Blue},
  urlcolor={blue},
  pdfcreator={LaTeX via pandoc}}

\author{}
\date{}

\begin{document}

\section{Four-Axis Computational Discovery in the eBL Cuneiform
Corpus}\label{four-axis-computational-discovery-in-the-ebl-cuneiform-corpus}

\subsection{Recovering Manuscript-Sibling False Negatives,
Reconstructing the Late-Mesopotamian Exorcist's Library, Discriminating
Scribal Lineages, and Restoring Damaged
Passages}\label{recovering-manuscript-sibling-false-negatives-reconstructing-the-late-mesopotamian-exorcists-library-discriminating-scribal-lineages-and-restoring-damaged-passages}

\begin{center}\rule{0.5\linewidth}{0.5pt}\end{center}

\textbf{Author}: Dane Brown --- Independent researcher, Tokyo, Japan ---
\texttt{dane@kairovault.com} \textbf{Code}:
\href{https://github.com/Hugegreencandle/cuneiform-mcp}{github.com/Hugegreencandle/cuneiform-mcp}
(commit \texttt{4976266}, v0.18.3) \textbf{License}: CC-BY 4.0
\textbf{Version}: arXiv preprint, 2026-05-17 (concurrent submission to
\emph{Cuneiform Digital Library Journal})

\begin{center}\rule{0.5\linewidth}{0.5pt}\end{center}

\subsection{Abstract}\label{abstract}

This paper documents a four-axis computational discovery pipeline
(\texttt{cuneiform-mcp} v0.16--v0.18.3) for the electronic Babylonian
Library transliteration corpus of approximately 36,500 tablets. The
pipeline operates on raw sign-token sequences without metadata and
produces four categories of finding: manuscript siblings missed by
lexical methods (fuzzy trigram-Jaccard with one-substitution tolerance
plus contiguous-run signaling), full manuscript-witness clusters of
canonical compositions (recursive breadth-first search via fuzzy
parallels), scribal-lineage candidates (orthographic-preference
clustering via log-likelihood-ratio signatures), and multi-sign lacuna
restorations (parallel-template alignment with bigram beam-search
fallback).

Four validated metrics result: a 55\% manuscript-sibling rescue rate on
candidates the strict 0.30 trigram-Jaccard methodology surfaces as
isolated; a 100+ tablet \emph{āšipūtu} (exorcist) canon cluster spanning
20 museum-collection prefixes recovered from a single seed tablet
(BM.77056); three reciprocal same-scribe pairs in 34 probed tablets with
empirically-validated discrimination from same-composition pairs; and
91.7\% top-1 precision with 100\% top-10 recall in synthetic-gap lacuna
restoration across 48 test cases.

The paper's primary methodological contribution is the demonstration
that exact lexical methods at the conventional 0.30 Jaccard threshold
systematically under-recover wide-transmission compositions, and that
the underlying \emph{āšipūtu} canon described in textual-evidence
consensus (Lenzi 2008; Geller 2010; Maul 1994) is recoverable
empirically from pure-orthographic clustering. A separate methodological
contribution emerges from two calibration audits demonstrating that
precision in cuneiform-discovery tooling is often calibration-limited
rather than signal-limited: a one-line scoring change lifted
lacuna-restoration top-1 precision from 22.9\% to 91.7\% without
altering the underlying recovery algorithm, and a threshold audit
converged the bi-orphan discovery surface from 167 candidates to 2
without altering any underlying method. A third contribution
demonstrates that the contiguous-run bonus is methodology-agnostic:
applied independently to fuzzy-trigram and exact-trigram methodologies,
both surface the same cross-subseries manuscript-section sibling pair
(K.5896 and K.2761, both Mīs pî, sharing continuous text passages with
K.2798, Bīt salāʾ mê) --- the calibration pattern itself transfers
across underlying algorithms.

\textbf{Keywords}: cuneiform; computational philology;
manuscript-witness reconstruction; sign-trigram methods; calibration
audit; \emph{āšipūtu} canon; Mīs pî; Bīt salāʾ mê.

\begin{center}\rule{0.5\linewidth}{0.5pt}\end{center}

\subsection{1. Introduction: The K.2798 ↔ Si.776 Case
Study}\label{introduction-the-k.2798-si.776-case-study}

K.2798 is a Neo-Assyrian tablet from Kuyunjik classified by the
electronic Babylonian Library (eBL) as
\texttt{CANONICAL/Magic/Purification/Bīt\ salāʾ\ mê}. Si.776 is a Sippar
tablet whose eBL description reads: ``Small tablet preserving the
beginning of the ritual tablet of Bīt šalāʾ mê. Written in Assyrian
script.'' They are confirmed manuscript siblings of the same canonical
purification ritual.

Yet lexical trigram-Jaccard rates the K.2798 ↔ Si.776 pair at 0.15 ---
below the conventional 0.30 discovery threshold (eBL's
\texttt{LineToVecRanker} default) --- because of two localized sign-form
variants at positions 4 and 5 that broke too many overlapping trigrams.
The first 12 of the first 14 sign tokens are identical between the two
tablets, but lexical methods at the standard threshold fail to surface
the relationship.

This pair motivates the entire methodology train documented here. Each
successive tool in the \texttt{cuneiform-mcp} pipeline performs a
specific function on this confirmed-sibling case:

\begin{enumerate}
\def\labelenumi{\arabic{enumi}.}
\tightlist
\item
  The bi-orphan anomaly surface (v0.16) flags K.2798 as a candidate
  worth probing.
\item
  The fuzzy parallel finder (v0.17) rescues the pair at fuzzy-Jaccard
  0.41 (a 2.67× lift over the exact-Jaccard score).
\item
  The cluster reconstructor (v0.17.1) reveals K.2798 as a peripheral
  witness of a wider Mīs pî + Bīt salāʾ mê cluster anchored at K.15325.
\item
  The scribal fingerprint (v0.18) confirms K.2798 and Si.776 are NOT
  same-scribe --- correctly distinguishing ``same composition'' from
  ``same scribe.''
\item
  The lacuna restorer (v0.18.1) correctly recovers synthetic gaps in
  K.2798 via the Si.776 template at 100\% top-10 recall.
\item
  The v0.18.2 thematic threshold audit lowers the bi-orphan threshold
  from 0.60 to 0.50 --- K.2798 ↔ Si.776 was thematic cosine 0.56, so the
  0.60 setting had been misclassifying confirmed sibling pairs as
  orphans.
\item
  The v0.18.3 run-bonus calibration ports a previously-shipped
  fuzzy-method scoring refinement to the exact-trigram tool, surfacing
  additional cross-subseries siblings (K.5896 and K.2761, both Mīs pî
  manuscripts) within K.2798's top-15.
\end{enumerate}

The K.2798 ↔ Si.776 case is the test case that validates each tool does
what it claims and motivates the calibration audits that produced this
paper's secondary methodological contribution.

\begin{center}\rule{0.5\linewidth}{0.5pt}\end{center}

\subsection{2. The Discovery Pipeline}\label{the-discovery-pipeline}

\subsubsection{2.1 Corpus and Filtering}\label{corpus-and-filtering}

The pipeline operates on the eBL \texttt{/api/fragments/all-signs}
endpoint dump (36,498 transliterated tablets, \textasciitilde33 MB),
filtered through a 20-record exclusion list of Asb.* colophon-template
prototype records that aggregate standardized Ashurbanipal-palace
colophon vocabulary across 200+ manuscripts each (Hunger 1968). After
exclusion and per-tool minimum-size filters, the working indices contain
28,665 tablets (thematic embeddings), 19,787 (lexical graph), 35,308
(fuzzy parallels), and 25,150 (scribal fingerprint).

\subsubsection{2.2 Lexical Layer}\label{lexical-layer}

Standard sign-trigram-Jaccard pairwise scoring with X-token filtering
(trigrams with ≥2 X positions skipped, per the 2026-05-14 X-FILTER
calibration). Threshold 0.30 (eBL's \texttt{LineToVecRanker} default).
The validation benchmark on 50 targets and 87 known siblings (combined
seed=42 + seed=137, N=267, 95\% CI {[}17\%, 28\%{]}) places exact
trigram-Jaccard at recall@15 of 22.5\%. Seventy-seven and a half percent
of known siblings score below top-15; most fall below 0.20 Jaccard.

\subsubsection{2.3 Thematic Layer}\label{thematic-layer}

Distributional semantic embeddings via Sahlgren (2005) Random Indexing:
300-dimensional, ±3 sign window, k=8 sparse nonzeros per index vector,
deterministic seed=42. Per-tablet vectors are IDF-weighted means of
sign-context vectors, mean-centered following Mu, Bhat, and Viswanath
(2018). Without mean-centering, random-pair cosine median was 0.97;
after, 0.00 with full spread from −0.70 to +0.80. Top-30 cosine
neighbors are precomputed per tablet.

\subsubsection{2.4 Anomaly Surface and Quality
Filters}\label{anomaly-surface-and-quality-filters}

A tablet with zero lexical neighbors at Jaccard ≥ 0.30 AND zero thematic
neighbors at cosine ≥ T (where T evolves through calibration) is a
\emph{bi-orphan}: isolated in both lexical and thematic spaces.

\begin{verbatim}
v0.16 threshold T = 0.60            → 167 bi-orphans corpus-wide
v0.18.2 threshold T = 0.50          → 11 bi-orphans
v0.18.2 + quality filters           → 2 bi-orphans (IM.49220 and K.3306)
\end{verbatim}

The threshold tightening from 0.60 to 0.50 is one of three calibration
audit results documented in §3. The K.2798 ↔ Si.776 confirmed-sibling
pair scores at thematic cosine 0.56 --- below 0.60 but above 0.50 --- so
the v0.16 setting was systematically capturing confirmed sibling pairs
as orphans.

Four quality filters apply default-ON for bi-orphan classification,
derived from a 2026-05-16 inspection of the top-15 bi-orphans: formulaic
(top-1 sign frequency \textgreater{} 12\%), refrain-heavy (max 3-gram
repetition \textgreater{} 3 in first 50 tokens), heavily-damaged
(X-ratio \textgreater{} 50\%), and provenance-cluster (\textgreater{}
80\% of top-30 thematic neighbors share a museum prefix).

\subsubsection{2.5 Fuzzy Lexical Layer}\label{fuzzy-lexical-layer}

Two trigrams \texttt{(a,\ b,\ c)} and
\texttt{(a\textquotesingle{},\ b\textquotesingle{},\ c\textquotesingle{})}
are fuzzy 1-substitution neighbors iff exactly 2 of 3 positions are
equal. For each tablet, three 2-of-3 inverted indices are built:
\texttt{(a,b)}, \texttt{(b,c)}, \texttt{(a,c)}. Fuzzy intersection is
the number of query trigrams with at least one fuzzy match in the
target.

The v0.18.2 calibration observed that bare fuzzy-Jaccard treated
scattered matches identically to contiguous runs. For K.2798 ↔ Si.776,
the 129 fuzzy matches form a 29-position contiguous run plus smaller
runs --- strong text-section sibling evidence the bare-Jaccard score did
not reflect. The calibration adds a run-bonus multiplier:

\begin{verbatim}
runFactor    = min(1, longest_run / sqrt(query_trigrams))
runBonus     = 0.5 × runFactor              (capped at +50% lift)
final_score  = fuzzy_jaccard × (1 + runBonus)
\end{verbatim}

\subsubsection{2.6 Cluster Reconstructor}\label{cluster-reconstructor}

Breadth-first search from a seed tablet. Each frontier node's top-K
fuzzy parallels with fuzzy-Jaccard ≥ threshold join the cluster and the
next frontier. Per-member topology preserved:
\texttt{(tablet\_id,\ depth,\ parent,\ fuzzy\_J\_to\_parent)}.
Termination occurs at depth-cap, size-cap, or frontier exhaustion.

\subsubsection{2.7 Scribal Fingerprint}\label{scribal-fingerprint}

Per-tablet \emph{scribal signature} = top-30 signs by log-likelihood
ratio of in-tablet vs.~corpus-baseline frequency:

\begin{verbatim}
LLR(s, T) = N_T(s) × log( P_T(s) / P_corpus(s) )
\end{verbatim}

filtered to signs with corpus frequency ≥ 5 and in-tablet count ≥ 2. Two
tablets with overlapping signatures share unusual orthographic
preferences --- variant-sign choices, logogram-vs-syllabic spelling
habits, sign-compound preferences. Comparison via signature cosine plus
signature Jaccard.

Because eBL transliterations normalize paleographic variation, this
measure captures spelling-preference fingerprint rather than handwriting
paleography in the strict sense.

\subsubsection{2.8 Lacuna Restorer}\label{lacuna-restorer}

For a damaged stretch of length k between known boundary signs: build
prefix and suffix context fingerprints from the ±W known signs adjacent
to the lacuna; scan the corpus for templates whose local sign sequence
contains BOTH a prefix-trigram AND a suffix-trigram within distance k ±
tolerance; extract intervening signs as candidate fills; score by:

\begin{verbatim}
score = sqrt(local_jaccard × bigram_coherence) × length_factor
\end{verbatim}

The length\_factor calibration is 1.0 for exact-length fills, 0.7 for
off-by-one, and 0.5 for off-by-two. This single-line refinement lifted
top-1 precision from 22.9\% (v0.18.0) to 91.7\% (v0.18.1) on the same
48-test benchmark without altering the underlying parallel-template
alignment algorithm.

If no parallel templates exist for a damaged passage, the methodology
falls back to bigram beam-search; this fallback is less reliable and
produces repetitive output on truly-novel material.

\begin{center}\rule{0.5\linewidth}{0.5pt}\end{center}

\subsection{3. Findings}\label{findings}

\subsubsection{\texorpdfstring{3.1 The BM.77056 Cluster: The
Late-Mesopotamian \emph{āšipūtu}
Library}{3.1 The BM.77056 Cluster: The Late-Mesopotamian āšipūtu Library}}\label{the-bm.77056-cluster-the-late-mesopotamian-ux101ux161ipux16btu-library}

Seeded at BM.77056 with \texttt{min\_fuzzy\_jaccard\ =\ 0.20}, the
cluster reconstructor at \texttt{max\_depth\ =\ 4},
\texttt{max\_size\ =\ 100} terminates at max-size-reached with 100
cluster members --- i.e., the underlying composition is wider than 100
manuscript witnesses.

The cluster spans 20 museum-collection prefixes: BM (31 members), K
(32), Sm (6), CBS (2), ND (2), N (1), IM (2), VAT (2), SU-1952 (1), UM
(1), Rm-IV (1), Rm-II (2), Ni (1), W (2), plus six accession-year ranges
in the British Museum's nineteenth-century acquisitions (1880, 1881,
1882, 1883, 1884) and one in modern accession (2023). Sixty-nine of 100
members are in a different prefix than the seed.

eBL metadata pulls on the top-fuzzy-Jaccard cluster members reveal
compositional content (table 1).

\textbf{Table 1}. eBL genre classifications for BM.77056 cluster top
members.

{\def\LTcaptype{none} % do not increment counter
\begin{longtable}[]{@{}
  >{\raggedright\arraybackslash}p{(\linewidth - 4\tabcolsep) * \real{0.3333}}
  >{\raggedright\arraybackslash}p{(\linewidth - 4\tabcolsep) * \real{0.3333}}
  >{\raggedright\arraybackslash}p{(\linewidth - 4\tabcolsep) * \real{0.3333}}@{}}
\toprule\noalign{}
\begin{minipage}[b]{\linewidth}\raggedright
Tablet
\end{minipage} & \begin{minipage}[b]{\linewidth}\raggedright
fuzzy-Jaccard
\end{minipage} & \begin{minipage}[b]{\linewidth}\raggedright
eBL genre
\end{minipage} \\
\midrule\noalign{}
\endhead
\bottomrule\noalign{}
\endlastfoot
K.15325, K.8994, K.11920 & --- &
\texttt{CANONICAL/Magic/Purification/Mīs\ pî} \\
K.2798, Si.776 & --- &
\texttt{CANONICAL/Magic/Purification/Bīt\ salāʾ\ mê} \\
BM.44414 & 0.23 & \texttt{CANONICAL/Magic/Exorcistic/Udugḫul} \\
BM.47782 & 0.21 & \texttt{CANONICAL/Literature/Hymns/Divine/Šuʾila}
(Namburbû? colophon) \\
BM.44035, BM.16745 & 0.22, 0.21 &
\texttt{CANONICAL/Magic/Anti-witchcraft} \\
BM.47910 & 0.24 & hymn to Nabû →
\texttt{CANONICAL/Literature/Hymns/Divine} \\
Sm.882 & 0.23 & \texttt{CANONICAL/Literature/Hymns/Divine} \\
BM.48903 (an isolate) & \textless0.20 &
\texttt{CANONICAL/Magic/Varia/Egalkura} \\
BM.35512 ↔ K.2581 (high-lift fuzzy pair) & 0.31 &
\texttt{CANONICAL/Technical/Medicine/Therapeutic/Recipe/Prescription\ (Šumma\ amēlu)} \\
\end{longtable}
}

This is the integrated \emph{āšipūtu} (exorcist) library --- the
canonical professional curriculum the exorcist's manual KAR 44 (Lenzi
2008) lists as comprising Mīs pî + Bīt salāʾ mê + Udugḫul + Šuʾila +
Namburbû + anti-witchcraft series + \emph{Šumma amēlu} medical
prescriptions + Egalkura palace-entering ritual. The corpus's
compositional unity, long recognized from textual and colophon evidence
(Lenzi 2008; Geller 2010; Maul 1994), is here demonstrated empirically
through pure-orthographic clustering at the sign-trigram level, with no
metadata input.

The exact-trigram-Jaccard methodology at the conventional 0.30 threshold
cannot recover this cluster as a connected component because peripheral
pair-similarities fall below threshold. Fuzzy 1-substitution Jaccard
plus BFS cluster reconstruction together produce the recovery.

\subsubsection{3.2 Manuscript-Sibling Rescue: 55\% of Filtered
Candidates}\label{manuscript-sibling-rescue-55-of-filtered-candidates}

For the 28 cleaned bi-orphans plus 3 high-lift fuzzy candidates (31
total) from the v0.17.0 surface, fuzzy parallels identify strong sibling
pairs (fuzzy-Jaccard ≥ 0.20) for 17 cases (55\%). Strongest rescued
pairs are listed in table 2.

\textbf{Table 2}. Top fuzzy-rescued sibling pairs with exact-Jaccard
lift.

{\def\LTcaptype{none} % do not increment counter
\begin{longtable}[]{@{}
  >{\raggedright\arraybackslash}p{(\linewidth - 8\tabcolsep) * \real{0.2000}}
  >{\raggedright\arraybackslash}p{(\linewidth - 8\tabcolsep) * \real{0.2000}}
  >{\raggedright\arraybackslash}p{(\linewidth - 8\tabcolsep) * \real{0.2000}}
  >{\raggedright\arraybackslash}p{(\linewidth - 8\tabcolsep) * \real{0.2000}}
  >{\raggedright\arraybackslash}p{(\linewidth - 8\tabcolsep) * \real{0.2000}}@{}}
\toprule\noalign{}
\begin{minipage}[b]{\linewidth}\raggedright
Seed → Sibling
\end{minipage} & \begin{minipage}[b]{\linewidth}\raggedright
fuzzy-J
\end{minipage} & \begin{minipage}[b]{\linewidth}\raggedright
exact-J
\end{minipage} & \begin{minipage}[b]{\linewidth}\raggedright
exact-J lift
\end{minipage} & \begin{minipage}[b]{\linewidth}\raggedright
Composition
\end{minipage} \\
\midrule\noalign{}
\endhead
\bottomrule\noalign{}
\endlastfoot
K.2798 → Si.776 & 0.4082 & 0.1528 & 2.67× & Bīt salāʾ mê \\
BM.34795 → BM.38295 & 0.3593 & 0.0115 & 31.2× & EAE Sîn (Reiner \&
Pingree 1998--2005) \\
BM.117666 → YBC.7455 & 0.3265 & --- & --- & Cross-collection \\
BM.35512 → K.2581 & 0.3096 & 0.0174 & 17.8× & \emph{Šumma amēlu}
medical \\
BM.45641 → BM.77056 & 0.2973 & 0.0180 & 16.5× & Joins the \emph{āšipūtu}
hub \\
\end{longtable}
}

The pattern of extreme exact-Jaccard lifts (17× and 31×) on confirmed
sibling pairs is methodologically significant: these are tablets where
exact methods score near-zero because sign-form variants do not
co-locate in trigram windows, even though compositional identity is
preserved.

\subsubsection{3.3 Cross-Subseries Discoveries Through Run-Bonus
Signaling}\label{cross-subseries-discoveries-through-run-bonus-signaling}

After the v0.18.2 fuzzy-parallels run-bonus calibration, the K.2798
ranking surfaces a previously-buried candidate: K.5896 at fuzzy-Jaccard
0.134 (below the typical 0.20 threshold) but contiguous run of 28
positions. The run-bonus lifts K.5896 to a strong-confidence discovery
candidate (final score 0.20). eBL probe confirms
\texttt{K.5896\ →\ CANONICAL/Magic/Purification/Mīs\ pî} (Neo-Assyrian
Kuyunjik).

The v0.18.3 port of the same calibration to
\texttt{find\_parallel\_text} (exact-trigram-Jaccard, the original v0.4
tool) independently confirms K.5896 --- and additionally promotes K.2761
(rank \#13 → \#9, run = 16 trigrams) into the top-10. eBL probe confirms
\texttt{K.2761\ →\ CANONICAL/Magic/Purification/Mīs\ pî} (Neo-Assyrian
Kuyunjik): another Mīs pî manuscript with a shared continuous text
passage to K.2798.

K.5896 and K.2761 are both Mīs pî manuscripts with continuous trigram
runs shared with K.2798 (Bīt salāʾ mê). The two compositions are sibling
subseries within the late-Mesopotamian purification-ritual canon --- and
the run-bonus signaling reveals empirically-real cross-subseries
manuscript-section siblings that bare aggregate-similarity scoring
buried in both methodologies. The independent convergence of v0.18.2
fuzzy and v0.18.3 exact methodologies on the same two Mīs pî tablets is
the strongest empirical validation that the calibration pattern itself
is the transferable contribution, not the specific algorithm.

\subsubsection{3.4 Scribal Validation: Reciprocal Pairs and Negative
Discrimination}\label{scribal-validation-reciprocal-pairs-and-negative-discrimination}

In the v0.18 scribal-fingerprint scale validation on 34 probed tablets,
three reciprocal same-scribe pairs emerged (table 3).

\textbf{Table 3}. Reciprocal same-scribe pairs (average signature
cosine).

{\def\LTcaptype{none} % do not increment counter
\begin{longtable}[]{@{}
  >{\raggedright\arraybackslash}p{(\linewidth - 6\tabcolsep) * \real{0.2500}}
  >{\raggedright\arraybackslash}p{(\linewidth - 6\tabcolsep) * \real{0.2500}}
  >{\raggedright\arraybackslash}p{(\linewidth - 6\tabcolsep) * \real{0.2500}}
  >{\raggedright\arraybackslash}p{(\linewidth - 6\tabcolsep) * \real{0.2500}}@{}}
\toprule\noalign{}
\begin{minipage}[b]{\linewidth}\raggedright
Pair
\end{minipage} & \begin{minipage}[b]{\linewidth}\raggedright
Composition
\end{minipage} & \begin{minipage}[b]{\linewidth}\raggedright
Average cosine
\end{minipage} & \begin{minipage}[b]{\linewidth}\raggedright
Combined evidence
\end{minipage} \\
\midrule\noalign{}
\endhead
\bottomrule\noalign{}
\endlastfoot
BM.77056 ↔ BM.74130 & \emph{āšipūtu} (Sippar) & 0.78 & fuzzy-J 0.48 +
reciprocal scribal --- probable physical same scribe \\
K.15325 ↔ K.8994 & Mīs pî (Kuyunjik) & 0.77 & fuzzy-J 0.49 + reciprocal
scribal at \#1/\#3 --- strongest pair \\
BM.35512 ↔ K.2581 & \emph{Šumma amēlu} medical & 0.59 & Cross-collection
same scribe candidate \\
\end{longtable}
}

The critical negative result: K.2798 ↔ Si.776 are NOT in each other's
same-scribe top-15. They are confirmed manuscript siblings (caught by
\texttt{find\_fuzzy\_parallels}) but the methodology correctly
identifies them as different scribes --- exactly as Assyriological
theory predicts for two scribes copying the same canonical text. This
negative result operationalizes the long-recognized distinction between
textual transmission and scribal lineage as two independent
computational objects: \texttt{find\_fuzzy\_parallels} answers ``what
composition?''; \texttt{find\_same\_scribe\_candidates} answers ``who
copied?''

A secondary discovery from the scribal pass concerns scribal-style
geography: Si.776 (Sippar provenance) has Kuyunjik-leaning scribal style
(K=5 of top-10 same-scribe candidates), corroborating eBL's catalog note
that it is ``written in Assyrian script.'' K.2798 (Kuyunjik provenance)
conversely has Babylonian-leaning scribal style (BM=5 of top-10). This
inverted pattern fits the Neo-Assyrian practice of importing Babylonian
scholarly material to Kuyunjik plus the reciprocal Neo-Babylonian
copying-from-Assyrian-models tradition (cf.~Beaulieu 2000).

\subsubsection{3.5 Lacuna Restoration: 92\% Top-1
Precision}\label{lacuna-restoration-92-top-1-precision}

Synthetic-gap stress test: 16 tablets × 3 lacuna sizes (3, 5, 7) = 48
test cases. Methodology: cut a clean window from a tablet's signs,
replace with X tokens, run the restorer, score against ground truth.

\textbf{Table 4}. Lacuna restorer performance by tablet and lacuna size,
before and after v0.18.1 length-factor calibration.

{\def\LTcaptype{none} % do not increment counter
\begin{longtable}[]{@{}llll@{}}
\toprule\noalign{}
Metric & v0.18.0 & v0.18.1 & Lift \\
\midrule\noalign{}
\endhead
\bottomrule\noalign{}
\endlastfoot
Top-1 exact-match precision & 22.9\% & \textbf{91.7\%} & +68.8 pts \\
Top-10 exact-match recall & 100.0\% & 100.0\% & preserved \\
Size 3 top-1 & 18.8\% & 87.5\% & +68.7 \\
Size 5 top-1 & 31.3\% & 93.8\% & +62.5 \\
Size 7 top-1 & 18.8\% & 93.8\% & +75.0 \\
\end{longtable}
}

The v0.18.0 → v0.18.1 calibration is the cleanest single illustration of
the broader methodological point: the underlying parallel-template
alignment had 100\% top-10 recall from v0.18.0 --- the signal was
present. Top-1 precision was 23\% because longer-fill templates
displaced exact-length matches to rank \#2--\#7. A single-line
length-factor multiplier lifted top-1 to 92\% on identical inputs.

The lacuna restorer also functions as a secondary
manuscript-sibling-recovery method: BM.46372 (a v0.17-filtered
``isolate'') is actually a join partner of BM.46338 per eBL's own
catalog, and the lacuna restorer recovers all three test windows at
top-1 because the join provides parallel templates.

\subsubsection{3.6 Discovery Surface Convergence to Two
Candidates}\label{discovery-surface-convergence-to-two-candidates}

{\def\LTcaptype{none} % do not increment counter
\begin{longtable}[]{@{}ll@{}}
\toprule\noalign{}
Stage & Bi-orphan count (≥100 signs) \\
\midrule\noalign{}
\endhead
\bottomrule\noalign{}
\endlastfoot
v0.16 raw & 42 \\
v0.17 (4 quality filters ON) & 28 \\
v0.17 + fuzzy rescue (55\%) & 11 still-isolated \\
v0.18.2 (threshold 0.50) & \textbf{2} \\
\end{longtable}
}

After all calibration, exactly two tablets remain: IM.49220 and K.3306.
These are the Class D ``numerical/metrological/zero-thematic-neighbors''
candidates from the original v0.16 inspection. The methodology converges
from a 36,500-tablet corpus to a 2-tablet genuinely-novel-composition
surface via empirical convergence --- three independent calibrations
(filter classes, threshold tightening, fuzzy rescue) all point to the
same final set.

The actual previously-unknown-composition yield is at most 2 of 42
candidates (approximately 5\%) --- well below the original v0.16
specification's target of 5/20. The bi-orphan methodology's real value
lies not in novelty discovery but in QA, cataloging-gap identification,
and manuscript-sibling false-negative recovery. Of the original 42
candidates: eight were already-published, known compositions (Ḫḫ IX-XI
lexical, eme.gi7 prayer XX, MUL.APIN astronomical, CT-published omens
and letters); one (BM.46372) was an explicit missed-join recorded in
eBL's own catalog. The methodology is QA tooling.

\begin{center}\rule{0.5\linewidth}{0.5pt}\end{center}

\subsection{4. The Calibration Audit
Methodology}\label{the-calibration-audit-methodology}

A separate methodological contribution emerges from two calibration
audits that demonstrated precision in cuneiform-discovery tooling is
often calibration-limited rather than signal-limited.

\subsubsection{4.1 The Lacuna Restorer Calibration (v0.18.0 →
v0.18.1)}\label{the-lacuna-restorer-calibration-v0.18.0-v0.18.1}

The underlying parallel-template alignment had 100\% top-10 recall from
v0.18.0 --- the signal was present. Top-1 precision was limited to 23\%
by a boundary-off-by-one bug in the ranking heuristic. A single-line
length-factor multiplier lifted top-1 to 92\% on identical inputs.

\subsubsection{4.2 The Bi-Orphan Surface Calibration (v0.16 →
v0.18.2)}\label{the-bi-orphan-surface-calibration-v0.16-v0.18.2}

Three concurrent fixes: (a) thematic threshold 0.60 → 0.50 (K.2798 ↔
Si.776 sibling pair was scoring at 0.56 and being misclassified as
orphan); (b) score-component rebalance (\texttt{sign\_count} was 77.7\%
of raw bi-orphan score; isolation strength was a minor modifier;
rebalanced to \texttt{isolation\ ×\ sqrt(sign\_count)}); (c) run-bonus
signaling for fuzzy-parallels (contiguous fuzzy-match runs were treated
identically to scattered noise; K.5896 Mīs pî discovery surfaced from
rank \textasciitilde9 to top-tier).

\subsubsection{4.3 The Methodology-Agnostic Run-Bonus
(v0.18.3)}\label{the-methodology-agnostic-run-bonus-v0.18.3}

Porting the v0.18.2 fuzzy run-bonus to \texttt{find\_parallel\_text}
(the v0.4 exact-trigram-Jaccard tool) independently surfaces the same
K.5896 cross-subseries discovery AND adds K.2761 to K.2798's top-15.
Both methodologies converge on the same Mīs pî ↔ Bīt salāʾ mê
manuscript-section-sibling pattern when given the same calibration
pattern. The calibration is methodology-agnostic.

\subsubsection{4.4 General Methodology}\label{general-methodology}

For each tool's scoring heuristic, the audit pattern: (1) decompose the
formula's component terms; (2) quantify per-component contribution on
real outputs (if one term dominates \textgreater70\%, investigate); (3)
test edge cases against known-positive pairs; (4) test with synthetic
ground truth; (5) ship empirically-validated fixes.

The v0.18.1 + v0.18.2 + v0.18.3 calibration trains collectively
establish that the signal is present; the ranking heuristic is the
lever. The independent convergence of v0.18.2 fuzzy and v0.18.3 exact
methodologies on the same K.5896 + K.2761 discovery --- driven by the
same run-bonus calibration pattern applied to two different underlying
algorithms --- is the strongest demonstration: the calibration pattern
itself is the transferable contribution.

\subsubsection{4.5 Audit Completion: Six Fixes, Two
No-Ops}\label{audit-completion-six-fixes-two-no-ops}

After two audit rounds, the cumulative tally is six fixes shipped
(lacuna length-factor; bi-orphan threshold tightening; bi-orphan score
rebalance; fuzzy run-bonus; find\_parallel\_text run-bonus; plus the
previously-shipped Mu \& Viswanath thematic mean-centering at index
build time) and two no-ops confirmed (thematic length-bias audit found
no bias; scribal threshold tuning found current calibration optimal).

The audit methodology correctly identified five of six tools with
ranking-stage biases that admit single-line or limited-scope fixes. Two
tools were confirmed clean --- themselves valuable findings: build-time
fixes (such as the v0.15 thematic-embedding mean-centering) do not
necessarily require ranking-stage calibration as well; and the
zero-triangle-clique result in scribal-fingerprint validation reflects a
real corpus structural property (scribes in the late-Mesopotamian record
predominantly worked in pairs of close orthographic kinship rather than
larger collaborative ateliers visible at LLR-signature level).

\begin{center}\rule{0.5\linewidth}{0.5pt}\end{center}

\subsection{5. Discussion}\label{discussion}

\subsubsection{5.1 What the Pipeline Does
Well}\label{what-the-pipeline-does-well}

\begin{enumerate}
\def\labelenumi{\arabic{enumi}.}
\tightlist
\item
  \textbf{Recovers wide-transmission manuscripts that exact methods
  miss.} Conventional 0.30 trigram-Jaccard is calibrated for nearby
  manuscripts; fuzzy 1-substitution plus recursive cluster expansion
  close the gap on peripheral witnesses.
\item
  \textbf{Surfaces compositional unity at the corpus level.} The
  BM.77056 result demonstrates that the late-Mesopotamian exorcist canon
  can be recovered as a single connected component by automated tooling,
  validating textual-evidence consensus with computational precision.
\item
  \textbf{Distinguishes orthogonal axes computationally.} The
  composition / scribe separation is validated cleanly by the K.2798 ↔
  Si.776 negative-result test.
\item
  \textbf{Achieves near-canonical precision in lacuna restoration.}
  91.7\% top-1 with 100\% top-10 recall on tablets with parallel
  templates.
\item
  \textbf{Converges the discovery surface via empirical calibration.}
  Three independent passes (quality filters, fuzzy rescue, threshold
  audit) collapse the bi-orphan surface to exactly the same 2
  candidates.
\end{enumerate}

\subsubsection{5.2 What the Pipeline Does Not
Do}\label{what-the-pipeline-does-not-do}

\begin{enumerate}
\def\labelenumi{\arabic{enumi}.}
\tightlist
\item
  \textbf{Does not discover previously-unknown compositions} at a rate
  above approximately 5\%. eBL editors have already curated
  genuinely-novel material; the bi-orphan surface mostly catches
  manuscript-sibling false negatives plus short-fragment material.
\item
  \textbf{Does not reconstruct full scribal ateliers.} Three reciprocal
  same-scribe pairs in 34 probed tablets but zero triangle cliques. The
  methodology surfaces strong pairs rather than entire collaborative
  groupings.
\item
  \textbf{Cannot handle joins that depend on metadata outside the
  indexed signs corpus} (the BM.46372 → BM.46338 case where the join is
  recorded in eBL's catalog but the linked tablet falls outside the
  signs corpus index).
\item
  \textbf{Beam-search fallback for truly-novel lacuna restoration is
  weak.} When no parallel templates exist, bigram beam-search collapses
  to repetitive high-frequency-sign output.
\end{enumerate}

\subsubsection{5.3 General Principles
Surfaced}\label{general-principles-surfaced}

Five general principles emerged from the v0.16 → v0.18.3 build sequence:

\begin{enumerate}
\def\labelenumi{\arabic{enumi}.}
\tightlist
\item
  Lexical thresholds calibrated for nearby manuscripts systematically
  under-recover wide-transmission canonical compositions.
\item
  Methodological precision in discovery tooling is often
  calibration-limited rather than signal-limited.
\item
  Threshold-calibration audits surface systematic mis-classification of
  genuine positives.
\item
  Score-component decomposition reveals single-axis dominance bugs.
\item
  The textual-transmission vs.~scribal-lineage distinction is
  operationalizable as two independent computational objects.
\end{enumerate}

\subsubsection{5.4 Implications for Future
Tooling}\label{implications-for-future-tooling}

The most natural next-step improvement is sign-form variant
normalization at tokenization time. The 2026-05-14 sign-variant
normalization experiment found zero rank-improvement from naive collapse
rules on an 87-sibling benchmark, but that experiment did not test
position-aware substitution at the trigram level --- which is what the
fuzzy method effectively does. A pre-pass that canonicalizes documented
sign-form-variant pairs into single tokens would yield equivalent
results at exact-Jaccard with lower compute cost.

Other natural extensions include compound-discovery tooling combining
all four signal axes; atelier reconstruction via recursive BFS in the
scribal-fingerprint graph; cross-corpus comparative tooling for Hebrew
Bible, Ugaritic, and Hittite parallels at the n-gram level; and
calibration audits on the lexical baseline
(\texttt{find\_parallel\_text} recall@15) and corpus-viz scoring.

\begin{center}\rule{0.5\linewidth}{0.5pt}\end{center}

\subsection{6. Reproducibility}\label{reproducibility}

All code is available at the cuneiform-mcp repository:
\href{https://github.com/Hugegreencandle/cuneiform-mcp}{github.com/Hugegreencandle/cuneiform-mcp}.
Commit \texttt{4976266} (v0.18.3) is the head documented here. The
repository is archived at Zenodo with DOI
\texttt{10.5281/zenodo.20250520}. All probes are deterministic (Random
Indexing seed = 42; BFS frontier order sorted by score within each
depth). Re-running reproduces the 100+ tablet BM.77056 cluster, the
17/31 fuzzy-rescue pairs, the 3 reciprocal same-scribe pairs, the 91.7\%
lacuna-restoration top-1 precision, the 2-tablet final bi-orphan
surface, and the K.5896 + K.2761 cross-subseries discoveries
identically.

Supplementary validation documents in the repository:

\begin{itemize}
\tightlist
\item
  \texttt{v0.16-bi-orphan-candidates.md} --- top-30 bi-orphan list
  (pre-filter)
\item
  \texttt{v0.16-bi-orphan-inspection.md} --- five-class false-positive
  analysis
\item
  \texttt{v0.17-validation-findings.md} --- 17/31 fuzzy-rescue results
\item
  \texttt{v0.17-ebl-metadata-probe.txt} --- raw eBL metadata pull for
  cluster identification
\item
  \texttt{v0.18-scribal-validation.md} --- three reciprocal pairs plus
  cross-axis discrimination
\item
  \texttt{v0.18-lacuna-stress-test.md} --- 48-test stress test, v0.18.0
  / v0.18.1 side-by-side
\item
  \texttt{v0.18.2-calibration-audit.md} --- three-fix audit details plus
  K.5896 discovery
\item
  \texttt{v0.18.3-calibration-round2.md} --- round-2 audit findings (two
  no-ops, one deferral resolved)
\item
  \texttt{v0.18.3-parallel-text-run-bonus.md} --- find\_parallel\_text
  calibration validation plus K.2761 discovery
\end{itemize}

\begin{center}\rule{0.5\linewidth}{0.5pt}\end{center}

\subsection{7. Conclusion}\label{conclusion}

This paper documents a four-axis computational discovery pipeline
operating on the eBL cuneiform corpus. The pipeline's primary
substantive contribution is the empirical reconstruction of the
late-Mesopotamian \emph{āšipūtu} (exorcist) library --- long recognized
from textual evidence --- as a 100+ tablet manuscript-witness cluster
recoverable from pure-orthographic clustering at the sign-trigram level.
The paper's primary methodological contribution is the demonstration
that exact lexical methods at conventional thresholds systematically
under-recover wide-transmission compositions, and that this
under-recovery is calibration-limited rather than signal-limited: a
one-line length-factor change lifted lacuna-restoration top-1 precision
from 22.9\% to 91.7\%, and a threshold audit converged the bi-orphan
discovery surface from 167 candidates to 2. A third contribution
demonstrates that the run-bonus calibration pattern is
methodology-agnostic, applying equivalently to fuzzy and exact
methodologies and surfacing the same K.5896 + K.2761 cross-subseries
discovery from both.

The pipeline's validated metrics --- 55\% sibling rescue rate, 100+
tablet cluster reconstruction, 3 reciprocal scribal pairs with
cross-axis discrimination, 91.7\% lacuna restoration top-1 --- establish
the methodology train as suitable for direct scholarly use on the eBL
corpus and as a transferable methodological pattern for other
ancient-text discovery tooling.

\begin{center}\rule{0.5\linewidth}{0.5pt}\end{center}

\subsection{Bibliography}\label{bibliography}

Beaulieu, P.-A. 2000. ``The Astronomers of the Esagil Temple in the
Fourth Century BC.'' In \emph{Assyriologica et Semitica: Festschrift für
Joachim Oelsner}. Münster: Ugarit-Verlag.

Geller, M. J. 2010. \emph{Ancient Babylonian Medicine: Theory and
Practice}. Chichester: Wiley-Blackwell.

Geller, M. J. 2016. \emph{Healing Magic and Evil Demons: Canonical
Udug-hul Incantations}. Berlin: De Gruyter. (Babylonisch-Assyrische
Medizin in Texten und Untersuchungen 8.)

Hunger, H. 1968. \emph{Babylonische und assyrische Kolophone}. Kevelaer
/ Neukirchen-Vluyn. (Alter Orient und Altes Testament 2.)

Lambert, W. G. 2007. \emph{Babylonian Oracle Questions}. Winona Lake,
IN: Eisenbrauns.

Lenzi, A. 2008. \emph{Secrecy and the Gods: Secret Knowledge in Ancient
Mesopotamia and Biblical Israel}. Helsinki: Neo-Assyrian Text Corpus
Project. (State Archives of Assyria Studies XIX.)

Maul, S. M. 1994. \emph{Zukunftsbewältigung: Eine Untersuchung
altorientalischen Denkens anhand der babylonisch-assyrischen Löserituale
(Namburbi)}. Mainz: P. von Zabern. (Baghdader Forschungen 18.)

Mu, J., Bhat, S., and Viswanath, P. 2018. ``All-but-the-Top: Simple and
Effective Postprocessing for Word Representations.'' In
\emph{Proceedings of the International Conference on Learning
Representations (ICLR)}.

Reiner, E., and Pingree, D. 1998--2005. \emph{Babylonian Planetary
Omens} 1--4. Leiden: Brill. (Cuneiform Monographs.)

Sahlgren, M. 2005. ``An Introduction to Random Indexing.'' Paper
presented at the Methods and Applications of Semantic Indexing Workshop
at the 7th International Conference on Terminology and Knowledge
Engineering (TKE).

Walker, C., and Dick, M. 2001. \emph{The Induction of the Cult Image in
Ancient Mesopotamia: The Mesopotamian Mīs Pî Ritual}. Helsinki:
Neo-Assyrian Text Corpus Project. (State Archives of Assyria Literary
Texts 1.)

\begin{center}\rule{0.5\linewidth}{0.5pt}\end{center}

\subsection{Appendix A: Build Sequence
(Supplementary)}\label{appendix-a-build-sequence-supplementary}

\begin{Shaded}
\begin{Highlighting}[]
\CommentTok{\# 1. Fetch eBL all{-}signs cache (\textasciitilde{}26 s, \textasciitilde{}33 MB)}
\ExtensionTok{node}\NormalTok{ dist/index.js }\AttributeTok{{-}{-}prefetch}

\CommentTok{\# 2. Build corpus{-}viz lexical graph (\textasciitilde{}7 min)}
\BuiltInTok{cd}\NormalTok{ \textasciitilde{}/Desktop/corpus{-}viz }\KeywordTok{\&\&} \ExtensionTok{node}\NormalTok{ build{-}graph.mjs}

\CommentTok{\# 3. Build v0.15 thematic embeddings (\textasciitilde{}4 min)}
\BuiltInTok{cd}\NormalTok{ \textasciitilde{}/Desktop/cuneiform{-}mcp }\KeywordTok{\&\&} \ExtensionTok{node}\NormalTok{ scripts/build{-}embeddings.mjs}

\CommentTok{\# 4. Build v0.16 anomaly{-}surface joined index (\textasciitilde{}10 s)}
\ExtensionTok{node}\NormalTok{ scripts/build{-}anomaly{-}index.mjs}

\CommentTok{\# 5. Reconstruct the BM.77056 cluster — single MCP call}
\BuiltInTok{echo} \StringTok{\textquotesingle{}\{"jsonrpc":"2.0","id":1,"method":"tools/call","params":\{"name":"reconstruct\_cluster","arguments":\{"seed\_tablet\_id":"BM.77056","max\_cluster\_size":100,"max\_depth":4\}\}\}\textquotesingle{}} \DataTypeTok{\textbackslash{}}
  \KeywordTok{|} \ExtensionTok{node}\NormalTok{ dist/index.js}
\end{Highlighting}
\end{Shaded}

\subsection{Appendix B: Validation
Benchmark}\label{appendix-b-validation-benchmark}

Validation benchmark for exact trigram-Jaccard recall@15 = 22.5\%
(combined seed=42 + seed=137, N=267, 95\% CI {[}17\%, 28\%{]}). Internal
documentation: \texttt{VALIDATION-2026-05-14.md} and
\texttt{TRIGRAM-EXPERIMENT-2026-05-14.md}.

\end{document}
